\documentclass[a4paper,12pt]{article}
\usepackage[utf8]{inputenc}
\usepackage[portuguese]{babel}
\usepackage{amsmath}
\usepackage{amsfonts}
\usepackage{listings}
\usepackage{xcolor}
\usepackage{geometry}
\usepackage{hyperref}
\hypersetup{hidelinks}
\usepackage{graphicx}

\geometry{margin=1in}

\definecolor{codegreen}{rgb}{0,0.6,0}
\definecolor{codegray}{rgb}{0.5,0.5,0.5}
\definecolor{codepurple}{rgb}{0.58,0,0.82}
\definecolor{backcolour}{rgb}{0.95,0.95,0.92}

\lstset{
    backgroundcolor=\color{backcolour},   
    commentstyle=\color{codegreen},
    keywordstyle=\color{magenta},
    numberstyle=\tiny\color{codegray},
    stringstyle=\color{codepurple},
    basicstyle=\ttfamily\footnotesize,
    breakatwhitespace=false,         
    breaklines=true,                 
    captionpos=b,                    
    keepspaces=true,                 
    numbers=left,                    
    numbersep=5pt,                  
    showspaces=false,                
    showstringspaces=false,
    showtabs=false,                  
    tabsize=2
}

\title{Lynceus Engine: Documentação Técnica de Extração Telemetria de Alta Fidelidade}
\author{Leonardo Herkenhoff \\ \small Instituto Federal do Espírito Santo (IFES) - PPComp}
\date{Abril de 2026}

\begin{document}

\maketitle

\begin{abstract}
Este documento detalha o funcionamento interno do motor Lynceus, uma infraestrutura baseada em eBPF/XDP para extração de telemetria de rede com precisão científica. O sistema implementa uma arquitetura híbrida de alto desempenho, combinando processamento em kernel-space para ingestão de pacotes e análise estatística online em user-space para derivação de mais de 495 features por fluxo.
\end{abstract}

\tableofcontents

\newpage

\section{Introdução}
O monitoramento de redes em alta velocidade exige soluções que minimizem o overhead de interrupções de hardware e cópias de memória. O Lynceus aborda este desafio utilizando o \textit{eXpress Data Path} (XDP), permitindo a dissecação de pacotes diretamente na camada de driver. Este documento serve como guia exaustivo para a compreensão da implementação v1.0.

\section{Arquitetura do Sistema}
O sistema é dividido em dois planos operacionais:
\begin{itemize}
    \item \textbf{Data Plane (Kernel):} Implementado em C/eBPF, responsável pela filtragem, decapsulamento recursivo e exportação via RingBuffers paralelos.
    \item \textbf{Control Plane (User):} Daemon multi-threaded que agrega eventos, aplica algoritmos estatísticos online e exporta os dados em formato CSV para pipelines de Machine Learning.
\end{itemize}

\section{Definições Compartilhadas (\texttt{lynceus.h})}
O arquivo \texttt{lynceus.h} estabelece a interface binária simétrica.

\subsection{Identidade de Fluxo (\texttt{flow\_id\_t})}
Utiliza-se uma estrutura compacta de 5-tuple (\textit{Protocol, IP Origem, IP Destino, Porta Origem, Porta Destino}). Para garantir compatibilidade dual-stack (IPv4/IPv6), os IPs são armazenados em arrays de 16 bytes.

\subsection{Vetor de Features (\texttt{flow\_record\_t})}
Estrutura que carrega metadados brutos do kernel, incluindo contadores de pacotes, bytes, flags TCP, tamanhos de janela e fragmentos de payload (\texttt{payload\_hint}) para análise de entropia.

\section{Data Plane em Kernel (\texttt{main.bpf.c})}
O programa XDP executa a lógica de decisão crítica.

\subsection{Dissecação de Protocolos}
O código realiza a descida recursiva na pilha de protocolos. 
\begin{enumerate}
    \item \textbf{L2/L3:} Identificação de EtherType e parsing de IPv4/IPv6.
    \item \textbf{Túneis:} Suporte nativo a GRE e VXLAN. Caso um cabeçalho UDP 4789 seja detectado, o motor realiza o salto para o payload interno.
    \item \textbf{L7 Discovery:} Parsers para DNS, SNMP, NTP e SSDP. Devido às restrições do \textit{BPF Verifier}, os loops de parsing L7 utilizam \texttt{\#pragma unroll} limitados a 16 iterações.
\end{enumerate}

\subsection{Gerenciamento de Mapas}
\begin{itemize}
    \item \textbf{\texttt{flow\_table}:} Tabela hash que mantém o estado efêmero do fluxo no kernel para evitar submissão redundante ao user-space em taxas de pacotes extremas.
    \item \textbf{\texttt{pkt\_ringbuf\_map}:} Um \textit{ARRAY\_OF\_MAPS} contendo RingBuffers dedicados por CPU (per-core), eliminando contenção de trava (\textit{lock contention}).
\end{itemize}

\section{Control Plane em User-Space (\texttt{loader.c})}
O daemon gerencia a complexidade estatística que o kernel não pode processar.

\subsection{Concorrência e Afinidade}
Cada CPU do sistema executa um Worker thread dedicado. A afinidade é configurada via \texttt{pthread\_setaffinity\_np} para garantir que os RingBuffers per-core sejam lidos localmente, maximizando o cache-locality.

\subsection{Matemática Estatística Online}
O Lynceus não armazena todos os pacotes. Ele utiliza algoritmos de passagem única.

\subsubsection{Algoritmo de Welford}
Para o cálculo de média ($\mu$), variância ($\sigma^2$), assimetria (\textit{skewness}) e curtose (\textit{kurtosis}) em uma única iteração, utiliza-se a atualização de momentos centrais:
\[
M_{k,n} = M_{k,n-1} + \dots \text{ (Fórmulas de momentos de 4ª ordem)}
\]
A variância é derivada como $\sigma^2 = \frac{M_2}{n-1}$.

\subsubsection{Estimativa de Quantis ($P^2$)}
Para medianas sem ordenação total, utiliza-se o algoritmo de Jain \& Chlamtac. Cinco marcadores rastreiam a densidade de probabilidade e se ajustam dinamicamente via interpolação parabólica.

\subsection{SPSC Queue e Exportação}
O sistema utiliza filas \textit{Single-Producer Single-Consumer} (SPSC) para desacoplar o processamento de rede da escrita em disco/stdout. Os Workers constroem as strings CSV em buffers de memória (\texttt{fmemopen}) e a thread de escrita (\textit{Writer}) realiza o drain assíncrono.

\section{Análise Crítica e Falhas Lógicas}
Como Parceiro de Pesquisa, aponto as seguintes limitações no código atual:
\begin{enumerate}
    \item \textbf{Janela de Observação L7:} O parser DNS está limitado a 16 labels. Pacotes maliciosos com profundidade maior evadirão a extração de metadados.
    \item \textbf{Fragmentação IP:} O motor assume que o cabeçalho L4 está presente no primeiro fragmento. Fragmentos subsequentes ou fora de ordem podem gerar 5-tuples incompletas.
    \item \textbf{Consumo de Memória:} A \texttt{flow\_table} em user-space é pré-alocada. Em cenários de ataques de exaustão de estado (SYN Flood massivo), a colisão de hash pode levar a perdas de telemetria se o scan de idle não for agressivo o suficiente.
\end{enumerate}

\section{Conclusão}
O Lynceus v1.0 representa uma base sólida para monitoramento avançado. A integração entre dissecadores de kernel e algoritmos estatísticos online permite uma visibilidade granular sem precedentes em fluxos de rede de alta velocidade.

\end{document}
